% Lösungen Einheit 2 – Nr. 30–39
\einheitenkopf[le2][2 Wachstumsfaktor und Tabelle]{Einheit 2 von 5 · Wachstumsfaktor und Wachstumstabelle}

\erg{30}{a) nimmt zu \quad b) nimmt ab \quad c) nimmt zu \quad d) nimmt ab \quad e) nimmt zu}
\erg{31}{a) 1,1 \quad b) 1,2 \quad c) 1,5 \quad d) 1,03 \quad e) 1,025 \quad f) 0,9 \quad g) 0,75 \quad h) 0,96 \quad i) 0,995 \quad j) 2}
\erg{32}{a) +4\,\% \quad b) +25\,\% \quad c) +8\,\% \quad d) +60\,\% \quad e) $-10$\,\% \quad f) $-15$\,\% \quad g) +1,5\,\% \quad h) $-0{,}5$\,\% \quad i) +200\,\%}
\erg{33}{a) 2 \quad b) 3 \quad c) 1,1 \quad d) 0,8 \quad e) ja, alle Quotienten sind 1,05; $q = 1{,}05$ \quad f) $q^2 = 484 : 400 = 1{,}21$, $q = 1{,}1$, also 10\,\% pro Jahr \quad g) $160 : 128 = 1{,}25$; $200 : 160 = 1{,}25$; $250 : 200 = 1{,}25$ – gleicher Faktor 1,25, also 25\,\% pro Stunde}
\erg{34}{a) 10; 20; 40 \quad b) 6; 18; 54 \quad c) 110; 121; 133,1 \quad d) 60; 90; 135}
\erg{35}{a) 250; 300; 360 (der Anfangswert steht bei $t = 0$) \quad b) 180; 259,2 \quad c) $800 \cdot 1{,}05^6 \approx 1072{,}08\,€$}
\erg{36}{a) $440 : 1{,}1 = 400$ \quad b) $972 : 128 = 7{,}59375 = 1{,}5^5$, also $t = 5$; Probe: $288 \cdot 1{,}5 \cdot 1{,}5 \cdot 1{,}5 = 972$ \quad c) 1800; 1350; 1012,5 \quad d) $34 \cdot 0{,}85 = 28{,}9$\,mg; $40 \cdot 0{,}85^4 \approx 20{,}88$, also $t = 4$\,h}
\erg{37}{a) Lena addiert jedes Jahr 20 statt mit 1,04 zu multiplizieren. Richtig: 500; 520; 540,80; 562,43 \quad b) Kai hat die Eins vergessen, der Faktor ist 1,05. Richtig: 200; 210; 220,50 \quad c) Jonas hat den Anfangswert schon einmal verzinst; bei $t = 0$ steht 800\,€. Richtig: 800\,€; 824\,€}
\erg{38}{a) Zum alten Wert (100\,\%) kommen 8\,\% dazu, zusammen 108\,\% = 1,08; mit 0,08 erhielte man nur den Zuwachs. \quad b) Nein: Die Quotienten $220 : 200$, $242 : 220$, $266{,}2 : 242$ sind alle 1,1. Bei exponentiellem Wachstum werden die Differenzen größer, der Quotient bleibt gleich.}
\erg{39}{a) 12; 15; 18,75; 23,44; 29,30; 36,62; 45,78 (in m², gerundet) \quad b) Woche 4 (29,30\,m² $>$ 25\,m²) \quad c) spätestens in Woche 5: dann sind es 36,62\,m², in Woche 6 schon 45,78\,m²}
